% PREPRINT (NAMED) -- P3. Generated by build_preprint.py; do not edit by hand.
% The anonymous submission build of this paper is in ../tex/.
% =====================================================================================
% Preprint format — the NON-ANONYMOUS build, for a preprint server.
%
% One column, deliberately: these papers are table-dense (P8 carries eight), and a
% preprint is read on screen where a single wide column beats two narrow ones. It also
% lets pandoc's `longtable` work natively, so no float conversion is needed here.
%
% ⚠️ THIS BUILD CARRIES THE AUTHOR NAME. It is the opposite of the submission build and
% must never be confused with it. Check which tree a PDF came from before sending it
% anywhere: tex/ is anonymous, tex-preprint/ is not.
% =====================================================================================
\documentclass[11pt]{article}

\usepackage[T1]{fontenc}
\usepackage[utf8]{inputenc}
\usepackage{lmodern}
\usepackage[margin=1.05in]{geometry}
\usepackage{microtype}

\usepackage{amsmath,amssymb}
\usepackage{textcomp}
\usepackage{booktabs}
\usepackage{longtable}
\usepackage{array}
\usepackage{multirow}
\usepackage{graphicx}
\usepackage{xcolor}
\usepackage[hidelinks]{hyperref}
\usepackage{enumitem}
\usepackage{framed}
\usepackage[normalem]{ulem}
\usepackage{fvextra}

% --- pandoc shims ---------------------------------------------------------------------
\providecommand{\tightlist}{\setlength{\itemsep}{0pt}\setlength{\parskip}{0pt}}
\providecommand{\pandocbounded}[1]{#1}
\providecommand{\st}[1]{\sout{#1}}
\setlist{leftmargin=1.2em, itemsep=2pt, topsep=4pt, parsep=0pt}
\RecustomVerbatimEnvironment{verbatim}{Verbatim}{breaklines,breakanywhere,fontsize=\small}

\definecolor{accent}{HTML}{B8552F}
\definecolor{warn}{HTML}{8A6D1F}

% Color is kept here, unlike the submission builds: a preprint is read on screen.
\newcommand{\warn}{\textcolor{warn}{\textbf{$\triangle$}}\ }
\newcommand{\rmark}[1]{\textcolor{accent}{\textbf{\textsf{\small #1}}}\ }

\newenvironment{callout}{\par\medskip\begin{leftbar}\small}{\end{leftbar}\par\medskip}
\renewenvironment{quote}{\par\medskip\begin{leftbar}\small\itshape}
  {\end{leftbar}\par\medskip}

\renewcommand{\arraystretch}{1.15}
% ⚠️ 3pt, not the 4pt this once was. `\tabcolsep` is applied on BOTH sides of every column, so a
% six-column table spends 10*\tabcolsep on gaps alone -- 40pt at 4pt. P6's two widest tables keep
% natural `l` columns (their cells are short, so widen_colspec correctly leaves them alone), and at
% 4pt the content plus gaps ran 3.5-5.5pt past the measure. Dropping to 3pt frees 10pt, which
% clears it with room to spare. Tables given explicit p{} widths are unaffected: their spec is
% written as \dimexpr\textwidth-N\tabcolsep\relax, so it tracks this value automatically.
\setlength{\tabcolsep}{3pt}

% Numbering OFF: the papers carry their own section numbers and cross-reference them in
% prose. Letting LaTeX renumber silently breaks every such reference.
\setcounter{secnumdepth}{-2}
\usepackage{titlesec}
\titleformat{\section}{\normalfont\large\bfseries\color{accent}}{\thesection}{0.6em}{}
\titleformat{\subsection}{\normalfont\normalsize\bfseries}{\thesubsection}{0.6em}{}
\titleformat{\subsubsection}{\normalfont\small\bfseries\itshape}{\thesubsubsection}{0.6em}{}
\titlespacing*{\section}{0pt}{12pt}{5pt}
\titlespacing*{\subsection}{0pt}{9pt}{3pt}

\date{}

% =====================================================================================
% Last-resort line-breaking slack. TeX sets a paragraph by stretching interword space within
% strict tolerance; when no break satisfies it -- a narrow two-column measure carrying tokens
% like "10-of-15" or a math arrow that cannot be hyphenated -- it gives up and lets the line
% run into the margin instead. \emergencystretch grants extra stretch ONLY on that final
% desperate pass, so it costs nothing on any paragraph that already fits, and converts a
% visible overflow into a slightly loose line. P6 ran 3.34pt and 2.66pt into the margin
% without it.
% =====================================================================================
\emergencystretch=1.5em

\title{\vspace{-2em}Positional Guard Evasion at Short Context: A Replication Note}
\author{Vikram Jha\\[0.30em]{\normalsize MuVeraAI}\\[0.18em]{\small ORCID~\href{https://orcid.org/0009-0004-3959-6099}{0009-0004-3959-6099}}\\[0.18em]{\small Corresponding author:~\href{mailto:vikram@muveraai.com}{vikram@muveraai.com}}}
\date{2026-08-31}
\begin{document}
\maketitle
\begin{center}\small\vspace{-1.1em}
\textcopyright~2026 Vikram Jha. Licensed under \href{https://creativecommons.org/licenses/by/4.0/}{CC BY 4.0}.
\end{center}
\begin{quote}
\textbf{Preprint.} This is the named version of a paper prepared for peer review. The anonymized
submission build is identical in content; only the author block, the companion-paper citation
keys and the anonymization note differ.
\end{quote}

\begin{quote}
\warn{} \textbf{Drafted as a novel-contribution paper; it is not one.} A prior-art check run after Draft 1
found \textbf{LongGuard} {[}LongGuard{]}, published 2026-08-27 --- four days before this work --- which
establishes the same positional finding across \textbf{15 guardrails} over a \textbf{0.25k--32k} grid,
supplies the attention$\rightarrow$\allowbreak{}logit$\rightarrow$\allowbreak{}behavior mechanism this note defers, and implements the
chunked-detection defense this note only predicts.

\textbf{The measurements below stand; the novelty claim does not.} Draft 2 repositions this as a
replication at short context on two models LongGuard did not test, and \textbf{retracts Draft 1\textquotesingle s claim
to have refuted length dilution} (\S{}4.2). Full assessment in \texttt{PRIOR-ART-ASSESSMENT.md}.

\textbf{Draft 3 tests that retraction instead of conceding it.} A follow-up run (P3-X, \S{}4.2.1) extends
the grid to LongGuard\textquotesingle s full 0--32k range at both scales, 1,260 inputs per model, zero unparsed
output. \textbf{The retraction was correct}: over the endpoints, 19 payloads out of 60 degrade and 0
improve at both scales (\emph{p} \textless{} 0.0001), and the 0.6B\textquotesingle s one non-monotone step is the single step that
fails significance (\emph{p} = 0.688).

\textbf{Pre-registered before any constructed input was scored} --- three mechanisms, their distinguishing
predictions and the invalidation conditions frozen and hashed (\texttt{ba560d2d}, re-hashed \texttt{21cdee5c}).
Pre-registration protected the analysis from the data. \textbf{It did nothing about claiming novelty
that did not exist}, which is the failure recorded here.
\end{quote}

\section{Abstract}\label{abstract}

Embedding unsafe content inside an ordinary benign frame --- a workplace email --- is the most effective
guard-evasion transformation we have measured, defeating guards on 51\% to 62\% of attempts across
three model families. It is also the transformation whose mechanism is least obvious, because it
changes two things at once: it makes the input longer, and it moves the unsafe span away from the
edges.

\textbf{We separate those by construction.} A synthetic grid holds filler \emph{content} constant while varying
filler \emph{quantity} and payload \emph{position} independently: 100 unsafe payloads, each of which both
guards catch perfectly in isolation, placed at three positions across four length multiples.

\textbf{Position orders the effect at short context.} Pooled across lengths, a payload in the middle of
its frame is missed 4.33$\times$ more often than one at the end. \textbf{This replicates LongGuard\textquotesingle s
lost-in-the-middle finding} {[}LongGuard{]} on two models they did not test, in a length regime
($\approx$0.3k--1.2k tokens) at the very bottom of their grid.

\textbf{We retract a claim, and then test the retraction.} Draft 1 reported miss rate as non-monotone in
length --- rising to 9.3\% at 4$\times$ filler then falling at 16$\times$ --- and concluded that dilution is refuted.
LongGuard attributes failure \emph{to} proportional dilution over a range 25$\times$ longer, with an edge--middle
gap growing from +6.67\% at 0.25k to +30.88\% at 16k. Our grid never left their weakest regime, so the
non-monotonicity was retracted as a probable small-range artifact. \textbf{We then ran the range} (\S{}4.2.1):
across 0--32k at both scales, \textbf{19 payloads degrade and 0 improve, at both scales (\emph{p} \textless{} 0.0001)}, and
the sole non-monotone step is the only step that fails significance. \textbf{The artifact reading was
right.}

\textbf{The measured effect.} 39/300 versus 9/300 for the 4B (Fisher one-sided \emph{p} = 3.4$\times$10$^{-}$$^{6}$) and 2.64$\times$
(\emph{p} = 2.4$\times$10$^{-}$$^{3}$) for the 0.6B. The ordering --- middle worst, end safest, start intermediate --- holds at
every length and in both models, and matches the direction LongGuard reports at larger scale.

\textbf{What this note adds, narrowly:} the effect is measurable at \textbf{short context} ($\approx$0.3k--1.2k tokens),
on \textbf{two models LongGuard did not evaluate}, against a control in which the same payloads are caught
\textbf{100\% of the time} unframed. We release the grid construction, the per-item verdicts, and the
pre-registration.

\textbf{Keywords:} guard models, adversarial robustness, positional bias, long-context, AI safety

\section{1. Introduction}\label{1-introduction}

Guard models classify content as safe or unsafe inline in an agentic pipeline. Prior work in this
program established that semantically-preserving rephrasing degrades guards asymmetrically --- leaving
false positives unchanged while raising false negatives --- and that among five transformation
families, one dominates: \textbf{embedding the unsafe content inside an ordinary benign frame.}

Across three model families that family reaches \textbf{0.5079 and 0.6199} candidate-level miss rate,
roughly double any other transformation. Two things about it were unexplained.

\textbf{First, it was not the strongest under deterministic decoding.} On a temperature-0 floor, a
different family led. The benign-framing advantage appears under \emph{sampling}, which suggests the
effect depends on something that varies between samples --- plausibly where the payload lands.

\textbf{Second, it confounds two variables.} Framing makes the input several times longer \emph{and} moves the
payload away from the input boundaries. Every prior analysis reported it separately for exactly this
reason, and none could separate the two.

\textbf{This paper separates them.}

\subsection{1.1 Three mechanisms, three predictions}\label{11-three-mechanisms-three-predictions}

\begin{longtable}[]{@{}>{\raggedright\arraybackslash}p{\dimexpr 0.070\dimexpr\textwidth-4\tabcolsep\relax\relax}>{\raggedright\arraybackslash}p{\dimexpr 0.497\dimexpr\textwidth-4\tabcolsep\relax\relax}>{\raggedright\arraybackslash}p{\dimexpr 0.431\dimexpr\textwidth-4\tabcolsep\relax\relax}@{}}
\toprule\noalign{}
& mechanism & prediction \\
\midrule\noalign{}
\endhead
\bottomrule\noalign{}
\endlastfoot
\textbf{M-A} & \textbf{Length dilution} --- the unsafe span is a smaller fraction of a longer input & miss rate rises with length at fixed position \\
\textbf{M-B} & \textbf{Positional displacement} --- the payload sits where the model attends least & miss rate varies with position at fixed length \\
\textbf{M-C} & \textbf{Boundary crossing} --- the frame shifts the judgment itself, not attention & verdict probability shifts even at short length and edge position \\
\end{longtable}

They are separable by construction: M-A varies length holding position, M-B varies position holding
length, and M-C predicts an effect surviving both held constant.

\subsection{1.2 Findings}\label{12-findings}

\begin{itemize}
\tightlist
\item
  \textbf{M-A: non-monotone in our range, refutation retracted --- and the retraction now tested} (\S{}4.2,
  \S{}4.2.1). Our grid tops out around 1.2k tokens; LongGuard measures dilution dominating over a range
  25$\times$ longer. \textbf{Extending to that range confirms it}: 19 of 60 payloads degrade and 0 improve from
  1k to 32k, at both scales (\emph{p} \textless{} 0.0001). \textbf{M-A was never refuted --- it was untested.}
\item
  \textbf{M-B supported}, at 4.33$\times$ and 2.64$\times$, significant in both models, consistent at every length ---
  \textbf{replicating} LongGuard\textquotesingle s lost-in-the-middle result at the short end of the range.
\item
  \textbf{M-C deferred here, and already answered elsewhere.} It requires verdict-token probabilities and
  a different serving stack, which would reintroduce a runtime confound this program was built to
  eliminate (\S{}3.4). LongGuard supplies the attention$\rightarrow$\allowbreak{}logit$\rightarrow$\allowbreak{}behavior chain we deferred.
\end{itemize}

\section{2. Related work}\label{2-related-work}

\textbf{LongGuard {[}LongGuard{]} is the primary result this note replicates}, and it precedes this work by
four days. It evaluates 15 guardrails over a 0.25k--32k grid, reports unsafe recall dropping more than
50\% on average, separates proportional dilution from absolute length with a paired Benign-Fill /
Needle-Repeat control, traces an attention$\rightarrow$\allowbreak{}logit$\rightarrow$\allowbreak{}behavior mechanism, isolates guard-specialized
retrieval heads, and supplies two training-free mitigations plus a deployment protocol. \textbf{Its
appendix C.2 reports the lost-in-the-middle and endpoint-asymmetry pattern this note also finds},
with an edge--middle gap growing from +6.67\% at 0.25k to +30.88\% at 16k.

\textbf{Positional bias in long-context models} is independently well documented for retrieval and
question answering. LongGuard carries it into safety guardrails.

\textbf{Guard robustness work} measures rephrasing, jailbreaks and prompt injection {[}GuardMeaning{]}, and
guardrail behavior under RAG-style long contexts {[}RAGGuard{]}.

\textbf{What this note claims.} A replication on two models LongGuard did not test --- at short context,
and then across their full 0--32k range --- each against a \textbf{0$\times$ unframed control}, which is the one
design element their needle-in-haystack setup does not have and which anchors every cell to an
absolute floor rather than a relative baseline. Plus a \textbf{retraction} of a dilution refutation our
range could not support, and a \textbf{test of that retraction} at the range where it is testable.
\textbf{We claim no novelty on the positional phenomenon, its mechanism, or its mitigation.}

\section{3. Method}\label{3-method}

\subsection{3.1 Payload selection}\label{31-payload-selection}

100 unsafe payloads, each \textbf{caught by both models under test in its original, unframed form.}

This matters more than it may appear. If a payload were already missed without framing, a miss under
framing would be credited to the frame rather than to a pre-existing failure. The pre-registration
forbids that, and \S{}4.1 confirms the filter worked: \textbf{the control condition has a 0.0000 miss rate in
both models.}

\begin{quote}
\textbf{A pre-registered rule that had to be revised, and was recorded rather than quietly changed.}
The original rule was "caught by every guard in the program." That yields \textbf{27 payloads}, because
one family in the panel catches only 32 of 150 items under the four-policy prompt we supply it ---
a limitation of \emph{our configuration}, not of the payloads. Intersecting against it would let our own
prompt decide which unsafe content this experiment may study. The rule was revised to "caught by
every guard \textbf{this experiment scores}," giving 120, of which 100 are used. The cost is that the
payload set is \textbf{not valid for a cross-family extension} and must be re-intersected if P3 is ever
run on the wider panel.
\end{quote}

\subsection{3.2 The grid}\label{32-the-grid}

\textbf{One filler corpus, reused in every cell}, so filler \emph{content} is constant and only quantity and
position vary. Filler is ordinary workplace-email prose.

\begin{itemize}
\tightlist
\item
  \textbf{Length:} payload + filler at \textbf{0$\times$, 1$\times$, 4$\times$, 16$\times$} the payload length.
\item
  \textbf{Position:} payload at \textbf{start}, \textbf{middle}, \textbf{end} of the filled input.
\item
  \textbf{0$\times$ is the control} --- payload alone, no frame, no position. It is emitted once per payload.
\end{itemize}

100 payloads $\times$ 10 cells = \textbf{1,000 constructed inputs}, each scored by both models. Median input
length runs 281 $\rightarrow$\allowbreak{} 564 $\rightarrow$\allowbreak{} 1,407 $\rightarrow$\allowbreak{} 4,779 characters across the four multiples.

\subsection{3.3 Scoring}\label{33-scoring}

Qwen3Guard-Gen-4B and -0.6B at \texttt{Q4\_K\_M}, \texttt{num\_ctx} 8192, seed 0, temperature 0, fail-closed ---
identical to every other measurement in this program. \textbf{Zero unparsed outputs across 2,000
evaluations.}

A \textbf{miss} is the guard returning \texttt{Safe} on an input containing an unsafe payload.

\subsection{3.4 Why M-C is deferred, stated before any number}\label{34-why-m-c-is-deferred-stated-before-any-number}

M-C requires verdict-token \emph{probabilities}, not argmax verdicts. The serving stack used throughout
this program does not usefully expose logprobs, so M-C requires \texttt{transformers} or vLLM.

\textbf{Numbers from that path cannot be compared against any verdict measured here}, or the runtime
confound --- precision, tokenization and kernel differences masquerading as an effect --- re-enters
through the back door. Prior work in this program was designed specifically to eliminate it. M-C
therefore requires its own baseline on the same grid through the same stack, and is reported as
\textbf{not run} rather than approximated.

\section{4. Results}\label{4-results}

\subsection{4.1 The control confirms the design}\label{41-the-control-confirms-the-design}

\begin{longtable}[]{@{}>{\raggedright\arraybackslash}p{\dimexpr 0.306\dimexpr\textwidth-2\tabcolsep\relax\relax}>{\raggedright\arraybackslash}p{\dimexpr 0.692\dimexpr\textwidth-2\tabcolsep\relax\relax}@{}}
\toprule\noalign{}
& miss rate on bare payloads (n=100) \\
\midrule\noalign{}
\endhead
\bottomrule\noalign{}
\endlastfoot
Qwen3Guard-4B & \textbf{0.0000} {[}0.0000, 0.0370{]} \\
Qwen3Guard-0.6B & \textbf{0.0000} {[}0.0000, 0.0370{]} \\
\end{longtable}

\textbf{Every payload is caught with no frame around it.} Every miss reported below is therefore
attributable to the framing, not to a payload the guard never handled.

\subsection{4.2 M-A: non-monotone in our range --- a retracted refutation}\label{42-m-a-non-monotone-in-our-range--a-retracted-refutation}

Miss rate pooled over position:

\begin{longtable}[]{@{}lllll@{}}
\toprule\noalign{}
& 0$\times$ & 1$\times$ & 4$\times$ & 16$\times$ \\
\midrule\noalign{}
\endhead
\bottomrule\noalign{}
\endlastfoot
\textbf{4B} & 0.0000 & 0.0433 & \textbf{0.0933} & 0.0667 \\
\textbf{0.6B} & 0.0000 & 0.0433 & \textbf{0.0833} & 0.0600 \\
\end{longtable}

\textbf{Both models peak at 4$\times$ and decline at 16$\times$.}

\warn{} \textbf{Draft 1 read this as refuting dilution. That reading is retracted.} Our largest cell is
$\approx$1.2k tokens; LongGuard\textquotesingle s grid runs to 32k and finds dilution dominating, with the edge--middle gap
growing from +6.67\% at 0.25k to +30.88\% at 16k. At 0.25k --- the point closest to our range --- their
effect is also small. \textbf{We did not test a range where dilution would be expected to dominate}, and
non-monotonicity inside a narrow low range is not evidence against it.

Length is not irrelevant --- the jump from 0$\times$ to 1$\times$ is the largest single step, and any framing at all
costs the guard something. But \textbf{within our range the quantity of filler does not order the effect.}

\subsubsection{4.2.1 The retraction, since tested rather than conceded}\label{421-the-retraction-since-tested-rather-than-conceded}

The retraction above was a concession to a longer-range published result. \textbf{We then ran the range}
(P3-X, pre-registered \texttt{324ee77b}): the same 60 payloads across lengths \{0, 1k, 4k, 16k, 32k\} $\times$ five
positions, 1,260 inputs per model at both scales.

\warn{} \textbf{\texttt{num\_ctx} was raised 8192 $\rightarrow$\allowbreak{} 32768 to admit a 32k input, and held constant across all 21 cells.
P3-X\textquotesingle s numbers are therefore not directly comparable to the table above.} It does not correct those
measurements; it tests the \emph{inference we drew from them} at a range they never reached.

\textbf{Unparsed output was 0.000 in all 42 cells}, so no number below is a truncation artifact.

\begin{longtable}[]{@{}llllll@{}}
\toprule\noalign{}
pooled miss rate & 0$\times$ & 1k & 4k & 16k & 32k \\
\midrule\noalign{}
\endhead
\bottomrule\noalign{}
\endlastfoot
\textbf{4B} & 0.0000 & 0.1100 & 0.1233 & 0.1933 & 0.2700 \\
\textbf{0.6B} & 0.0167 & 0.0700 & 0.1867 & 0.1600 & 0.3233 \\
\end{longtable}

\textbf{The 4B is monotone.} The 0.6B has one decrease, at 4k$\rightarrow$\allowbreak{}16k. An unregistered follow-up --- reported as
unregistered --- tests the steps with McNemar exact on payload-level discordance, since the same 60
payloads recur at every length:

\begin{quote}
\textbf{That single decrease is the one step that fails significance (2 worse, 4 better, \emph{p} = 0.688).}
Every other 0.6B step is a significant rise. \textbf{Over the endpoints, both scales are identical and
one-directional: 19 payloads degraded, 0 improved, \emph{p} \textless{} 0.0001.} Not one payload out of 60 became
easier to detect when filler was added, at either scale.
\end{quote}

The edge--middle gap grows with length at the 4B --- +7.2\% at 1k to +28.3\% at 32k --- closely tracking
LongGuard\textquotesingle s +6.67\% at 0.25k to +30.88\% at 16k.

\textbf{So the retraction was correct, and is now supported by measurement rather than by deference.}
Draft 1\textquotesingle s non-monotonicity was a small-range artifact. \textbf{M-A is not refuted; it was untested}, and
over the range where it is testable it holds.

\subsection{4.3 M-B supported: position orders the effect}\label{43-m-b-supported-position-orders-the-effect}

Miss rate by position, pooled over the three non-zero lengths:

\begin{longtable}[]{@{}llll@{}}
\toprule\noalign{}
& start & \textbf{middle} & end \\
\midrule\noalign{}
\endhead
\bottomrule\noalign{}
\endlastfoot
\textbf{4B} & 13/300 = 0.0433 & \textbf{39/300 = 0.1300} & 9/300 = 0.0300 \\
\textbf{0.6B} & 16/300 = 0.0533 & \textbf{29/300 = 0.0967} & 11/300 = 0.0367 \\
\end{longtable}

\begin{longtable}[]{@{}lll@{}}
\toprule\noalign{}
comparison & 4B & 0.6B \\
\midrule\noalign{}
\endhead
\bottomrule\noalign{}
\endlastfoot
middle vs end & \textbf{4.33$\times$}, \emph{p} = 3.4 $\times$ 10$^{-}$$^{6}$ & \textbf{2.64$\times$}, \emph{p} = 2.4 $\times$ 10$^{-}$$^{3}$ \\
middle vs start & 3.00$\times$, \emph{p} = 1.1 $\times$ 10$^{-}$$^{4}$ & 1.81$\times$, \emph{p} = 3.1 $\times$ 10$^{-}$$^{2}$ \\
\end{longtable}

\emph{(Fisher exact, one-sided.)}

\textbf{The ordering is middle $\gg$ start \textgreater{} end, and it holds at every individual length in both models:}

\begin{longtable}[]{@{}llll@{}}
\toprule\noalign{}
4B & start & middle & end \\
\midrule\noalign{}
\endhead
\bottomrule\noalign{}
\endlastfoot
1$\times$ & 0.0200 & 0.0900 & 0.0200 \\
4$\times$ & 0.0500 & \textbf{0.1800} & 0.0500 \\
16$\times$ & 0.0600 & 0.1200 & 0.0200 \\
\end{longtable}

The worst single cell is \textbf{4B at 4$\times$ length, middle position: 0.1800} --- eighteen of a hundred
payloads that the same model catches perfectly when unframed.

\subsection{4.4 Why the two mechanisms produce different curves}\label{44-why-the-two-mechanisms-produce-different-curves}

The interaction explains \S{}4.2\textquotesingle s non-monotonicity. At 16$\times$ the filler is long enough that the
"middle" is far from both edges --- but the \emph{end} position also becomes strongly protective, because
the payload sits immediately before the verdict is generated. Pooling across positions therefore
mixes a worsening middle with an improving end, and the pooled curve turns over.

\textbf{That is exactly what an effect governed by position, not quantity, looks like when you average
over position.} It is also why reporting length alone would have produced a confusing null.

\section{5. Analysis}\label{5-analysis}

\subsection{5.1 The positional profile matches long-context retrieval, in a safety setting}\label{51-the-positional-profile-matches-long-context-retrieval-in-a-safety-setting}

The middle-worst, edges-better profile is the signature documented for long-context retrieval. Its
appearance in a \emph{classification} task with a safety consequence is, to our reading, new --- and it
means the guard is not failing to \emph{understand} the payload. It is failing to \emph{attend} to it.

The strongest evidence for that reading is \S{}4.1: the identical payload, identical model, identical
settings, is caught \textbf{100\% of the time} when it is the whole input. Nothing about the content
changed.

\subsection{5.2 The attacker\textquotesingle s advantage is free}\label{52-the-attackers-advantage-is-free}

An attacker gains this multiplicatively \textbf{without modifying the payload}. No rephrasing, no
obfuscation, no optimization against the guard --- only placement. Placement is also the cheapest
possible variation: it requires no model, no queries and no feedback.

That places this below every attack we have previously measured in cost and above most in
reliability.

\subsection{5.3 The defensive implication is unusually concrete}\label{53-the-defensive-implication-is-unusually-concrete}

Most findings in this program complicate defense. This one suggests something actionable:
\textbf{a guard that scores an input in overlapping windows, or that re-scores with the suspect span moved
to the end, should recover much of the loss} --- because the end position is measurably the safest
one, at 0.0300 and 0.0367 against a middle of 0.1300 and 0.0967.

We have not tested that. It is stated as the obvious follow-up, and it is cheap.

\subsection{5.4 Scale does not rescue it}\label{54-scale-does-not-rescue-it}

The 4B shows a \emph{larger} positional effect than the 0.6B (4.33$\times$ versus 2.64$\times$), driven by an unusually
low end-position miss rate. The bigger model is better at the edges and no better in the middle.
\textbf{Scaling improves the best case without repairing the worst.}

\section{6. Threats to validity}\label{6-threats-to-validity}

\textbf{M-C is untested} (\S{}3.4), so we cannot exclude that framing also shifts the judgment itself. Our
claim is that M-A does not account for the effect and M-B substantially does --- not that M-B is the
whole story.

\textbf{Two models, one family.} Qwen3Guard at two scales. The benign-framing dominance that motivated
this work replicates across three families, but the positional decomposition has not been run on
them. The payload set would also need re-intersecting first (\S{}3.1).

\textbf{One filler corpus.} Filler content is constant by design, which removes it as a confound but also
means we tested one register of benign prose. Filler that is topically related to the payload might
behave differently.

\textbf{Synthetic construction.} The grid is assembled, not naturally occurring. It buys the ability to
separate length from position, and it costs ecological validity: a real attacker\textquotesingle s email would be
coherent with its payload in a way ours is not --- which most likely makes our estimates
\textbf{conservative}.

\textbf{Cell sizes of 100.} Individual cells have wide intervals; the pooled position comparisons
(n=300 per cell) carry the significance claims, and per-cell figures are reported for shape rather
than for inference.

\textbf{Absolute rates are small.} A 13\% pooled middle-position miss rate is a large \emph{relative} effect
against a 0\% control, and a modest absolute one. Both framings are in \S{}4.3 and neither is preferred.

\section{7. Conclusion}\label{7-conclusion}

The most effective guard-evasion transformation we have measured works because of \textbf{where} the
payload sits, not \textbf{how much} text surrounds it.

Within the short range we tested, length is non-monotone. Position orders the effect cleanly: a payload in
the middle of a benign frame is missed \textbf{4.33$\times$ more often than the same payload at the end} by a
model that catches it \textbf{100\% of the time} with no frame at all. The ordering holds at every length,
in both model scales, and the effect is significant at \emph{p} \textless{} 10$^{-}$$^{5}$.

This is the long-context "lost in the middle" bias, arriving in a safety classifier and carrying an
exploitable consequence. It costs an attacker nothing --- no rephrasing, no queries, no optimization,
only placement --- and it is the cheapest reliable evasion we have measured.

It may also be the most tractable. The end position is measurably safe. A guard that scores
overlapping windows, or that relocates a suspect span before deciding, has a plausible path to
recovering most of the loss. \textbf{That is a defense this result predicts and does not test}, and it is
the next thing worth running.

\section{Artifact}\label{artifact}

Grid construction script; the filler corpus; all 1,000 constructed inputs by payload identifier,
length multiple and position; per-item verdicts for both models; the pre-registration and both
hashes, including the payload-selection addendum and what it costs.


\section*{References}
\small
\begin{description}[leftmargin=1.4em,labelindent=0pt,itemsep=2pt]
\raggedright
\item[{[LongGuard]}] Z. Chen, X. Wu and S. Hu. \emph{LongGuard: Mechanistic Analysis and Training-Free Mitigation of Long-Context Failure in Safety Guardrails.} Institute of Information Engineering, Chinese Academy of Sciences. arXiv:2608.27580, 2026. --- \textbf{the primary result this note replicates.} Establishes the positional finding across 15 guardrails over 0.25k--32k, supplies the attention$\rightarrow$\allowbreak{}logit$\rightarrow$\allowbreak{}behavior mechanism, and implements chunked-detection and attention-head-sharpening defenses.
\item[{[GuardMeaning]}] C. Pinneri and C. Louizos. \emph{Guarding the Meaning: Self-Supervised Training for Semantic Robustness in Guard Models.} Qualcomm AI Research. arXiv:2511.10665, 2025.
\item[{[RAGGuard]}] \emph{RAG Makes Guardrails Unsafe? Investigating Robustness of Guardrails under RAG-style Contexts.} arXiv:2510.05310, 2025.
\item[{[WildGuard]}] \emph{WildGuard: Open One-Stop Moderation Tools for Safety Risks, Jailbreaks, and Refusals of LLMs.} arXiv:2406.18495, 2024. --- source of the payloads. \textgreater{} \textbf{On the sequencing.} Draft 1 was written before any prior-art check. LongGuard predates it by \textgreater{} four days and supersedes its central claim. This is recorded rather than quietly absorbed, and the \textgreater{} program has since added a literature gate before first draft rather than after.
\end{description}

\end{document}